\documentclass[12pt,a4paper]{article}
\usepackage[margin=25mm, headheight=15pt, headsep=10pt]{geometry}
\usepackage{fontspec}
\usepackage[table]{xcolor} % Note: table option is required for rowcolors
\usepackage{titlesec}
\usepackage{tocloft}
\usepackage{fancyhdr}
\usepackage{booktabs}
\usepackage{tcolorbox}
\tcbuselibrary{skins, breakable}
\usepackage[colorlinks=true, linkcolor=emerald, urlcolor=emerald, bookmarks=true]{hyperref}

\definecolor{emerald}{RGB}{0,102,51}
\definecolor{darkred}{RGB}{139,0,0}
\definecolor{lightgray}{RGB}{245,245,245}

\usepackage[nil]{babel}
\babelprovide[import, main]{bengali}
\babelprovide[import]{arabic}
\babelfont{rm}[Renderer=HarfBuzz,Scale=1.0]{Noto Serif Bengali}
\babelfont{sf}[Renderer=HarfBuzz]{Noto Sans Bengali}
\babelfont{tt}[Renderer=HarfBuzz]{Noto Sans Bengali}
\babelfont[arabic]{rm}[Renderer=HarfBuzz,Scale=1.1]{Noto Naskh Arabic}

% Section styling
\titleformat{\section}{\Large\bfseries\color{emerald}}{\thesection.}{0.5em}{}
\titleformat{\subsection}{\large\bfseries\color{emerald!85!black}}{\thesubsection.}{0.5em}{}
\titleformat{\subsubsection}{\normalsize\bfseries\color{emerald!70!black}}{\thesubsubsection.}{0.5em}{}
\titlespacing*{\section}{0pt}{18pt}{8pt}

% Fancy Header/Footer setup
\pagestyle{fancy}
\fancyhf{} % clear all header and footer fields
\fancyhead[L]{\color{emerald}\small\textbf{\nouppercase{\leftmark}}}
\fancyfoot[L]{\scriptsize\color{black!50}{\lat{AI}}-সংকলিত, আলেম-যাচাইকৃত নয়}
\fancyfoot[C]{\color{emerald!70!black}\thepage}
\renewcommand{\headrulewidth}{0.4pt}
\renewcommand{\headrule}{\hbox to\headwidth{\color{emerald}\leaders\hrule height \headrulewidth\hfill}}
\renewcommand{\footrulewidth}{0pt}

% Custom boxes
% 1. Ayat/Hadith Box
\newtcolorbox{ayatbox}[1][]{
    enhanced, breakable, 
    colback=emerald!5, 
    colframe=emerald, 
    boxrule=0.5pt, 
    arc=3pt, 
    left=10pt, right=10pt, top=10pt, bottom=10pt,
    #1
}

% 2. Quote/Translation Box
\newtcolorbox{quotebox}[1][]{
    enhanced, breakable,
    frame hidden,
    colback=lightgray,
    borderline west={3pt}{0pt}{emerald},
    left=15pt, right=10pt, top=10pt, bottom=10pt,
    sharp corners,
    #1
}

% 3. AI-generated notice box
\newtcolorbox{warnbox}[1][]{
    enhanced, breakable,
    colback=red!4,
    colframe=darkred,
    boxrule=0.8pt,
    arc=3pt,
    left=10pt, right=10pt, top=10pt, bottom=10pt,
    #1
}

% Commands
\newcommand{\arab}[1]{\foreignlanguage{arabic}{#1}}
\newcommand{\kib}[1]{\textcolor{emerald}{\textbf{#1}}}

% Latin (English) fallback: Noto Serif Bengali has NO Latin glyphs
\newfontfamily\latfont[Renderer=HarfBuzz]{Noto Serif}
\newcommand{\lat}[1]{{\latfont #1}}

\setlength{\parskip}{5pt}
\setlength{\parindent}{0pt}

% Fix for missing bullet in Noto Serif Bengali
\renewcommand{\labelitemi}{$\bullet$}



\begin{document}
\begin{center}{\Large\bfseries\color{emerald} পার্ট ৭ক — নামাজ ভঙ্গকারী ও নিষিদ্ধ বিষয়}\end{center}
\vspace{1em}
\section*{পার্ট ৭ক — নামাজ ভঙ্গকারী ও নিষিদ্ধ বিষয়}

\begin{warnbox}
 \textbf{গুরুত্বপূর্ণ নোটিশ (\lat{Important} \lat{Notice}):} এই কন্টেন্টটি \lat{AI} (কৃত্রিম বুদ্ধিমত্তা) দিয়ে প্রস্তুত ও সংকলিত। \lat{AI} কঠোর সূত্র-মান নীতি মেনে শুধু নির্ভরযোগ্য উৎস থেকে তথ্য সংগ্রহ করেছে — কেবল সহীহ/হাসান হাদিস ও সহীহ সনদের আসার রাখা হয়েছে, দুর্বল সনদ বাদ দেওয়া হয়েছে। তবে এটি কোনো আলেম বা মুহাদ্দিস কর্তৃক যাচাইকৃত নয়।
\end{warnbox}

\begin{quote}
\textbf{সম্পর্কিত ফাইল:} পার্ট ৩ (শ্রেণিবিভাগ), পার্ট ৫খ (সাহু সিজদা), পার্ট ৭খ (ভুল-ও-প্রতিকার)।
\end{quote}

\begin{quote}
\textbf{যে প্রশ্নের উত্তর এই অংশে:} নামাজে কী কী করলে নামাজ \textbf{ভেঙে যায়} (মুফসিদ)? কোন সময়ে নামাজ \textbf{নিষিদ্ধ}? নামাজে কথাবার্তা, খাওয়া, হাসি-কান্না — কী বিধান?
\end{quote}

\medskip\hrule\medskip

\subsection*{১. নামাজ ভঙ্গকারী (\lat{Mufsidaat}) — প্রধান তালিকা}

\textbf{নামাজ ভঙ্গ করে} (সব মাযহাবের মোটা ঐকমত্য, বিস্তারিত মাযহাবগত পার্থক্য নিচে):

\begin{enumerate}
\item \textbf{ইচ্ছাকৃত কথা বলা} — নামাজে কোনো কথাই জায়েয নয়।
\end{enumerate}
\begin{itemize}
\item \textbf{সহীহ হাদিস (সহীহ বুখারি ১২১৭; সহীহ মুসলিম ৫৩৭):} একবার সাহাবিরা নামাজে কথা বলছিলেন; নবীজি (সা.) বললেন — \textbf{\lat{“}নামাজ তাসবিহ, তাকবির ও কুরআন তিলাওয়াতের জন্যই; তাতে মানুষের কথা বলার জায়গা নেই।\lat{”}}
\item \textbf{মতভেদ:} ভুলে কথা বললে — হানাফি: নামাজ ভাঙে (কিন্তু সাহু সিজদা নয়); জমহুর: নামাজ ভাঙে (ইচ্ছা-ভুল একই)। মালেকি-হাম্বলি: ভুলে কথা বললে সাহু সিজদা দিয়ে ঠিক হওয়া যায়।
\end{itemize}
\begin{enumerate}
\item \textbf{খাওয়া-পান করা} — ইচ্ছাকৃত খাওয়া-পান নামাজ ভাঙে (সব মাযহাব)।
\item \textbf{বড় আন্দোলন} (কাজস-কাসীরা): নামাজের মতো স্থিরতা নষ্ট করা, টানা-বারবার চলা — যেমন দুই-তিন ধাপ পর পর এগোনো — নামাজ ভাঙে (মাযহাবভেদে মানদণ্ড: হানাফি "অনেক" — জমহুর "সংখ্যা/প্রকৃতিতে" ব্যাপক)।
\item \textbf{বড় অপবিত্রতা (হাদাস) হওয়া} — প্রস্রাব-পায়খানা-বায়ু বের হওয়া; নামাজ ভেঙে যায় (সব মাযহাব)। বায়ুর ক্ষেত্রে কিছু মাযহাবের সামান্য পার্থক্য (নিষ্ক্রিয় শিথিলতা-দর্শন)।
\item \textbf{অন্তর ফিরে যাওয়া (খেয়াল-বিচ্যুতি)} — সাধারণত নামাজ ভাঙে না, তবে খুশু নষ্ট হয়; যদি এমন অবস্থায় ঘোরাঘুরি বড় হয় — ভাঙে।
\item \textbf{জানামতে সতর খুলে যাওয়া} — পুরুষের নাভি-হাঁটু বা নারীর সতর উন্মুক্ত হলে, ইচ্ছাকৃতভাবে ঢাকলে না — নামাজ ভাঙে; যদি অজান্তে খুলে গিয়ে তৎক্ষণাৎ ঢাকেন — ভাঙে না (মাযহাবগত শিথিলতা)।
\end{enumerate}
\medskip\hrule\medskip

\subsection*{২. নামাজে নিষিদ্ধ (নামাজ ভাঙে না, তবে গুনাহ/মাকরূহ)}

\begin{itemize}
\item \textbf{মোবাইল-ফোন বেজে ওঠা:} নামাজ ভাঙে না, তবে \textbf{মাকরূহ} (দৃষ্টি-একাগ্রতা নষ্ট); অ্যাপ-রিডিং/চলন্ত মোবাইল সেজদার দিকে — খুশু নষ্ট করে, সুতরাং সেটা পরিহার (আধুনিক ফতোয়া)। \textbf{মোবাইল নীরব রেখে নামাজে যাওয়া} — মুস্তাহাব/আদব।
\item \textbf{ঘোরাঘুরি, চুল-কাপড় বোলানো, তাকানো} — মাকরূহ (পার্ট ৩-এর ৬ নম্বর অংশ)।
\item \textbf{খেলনা-ছবি সামনে থাকা:} যদি খুশু নষ্ট হয়, তাহলে তা সরিয়ে রাখা; নবীজি (সা.) নিজে ছবি-যুক্ত কাপড়ে নামাজ অপছন্দ করেছেন (সহীহ বুখারি ৩৭৪ — আয়েশা (রাঃ)-এর ছবি-চাদরের ঘটনা)।
\item \textbf{পায়খানা-প্রস্রাবের চাপে নামাজ:} হাদিসে \lat{“}খাবার হাজির থাকলে ও পেশাব-পায়খানার চাপ থাকলে নামাজ পড়ো না\lat{”} (সহীহ মুসলিম ৫৬০) — খুশু নষ্টের কারণে মাকরূহ; তবে ওয়াক্ত শেষ হয়ে যাওয়ার আশঙ্কায় পড়া ওয়াজিব।
\end{itemize}
\medskip\hrule\medskip

\subsection*{৩. নিষিদ্ধ সময় (নামাজ না পড়ার সময়)}

\textbf{সহীহ হাদিস (সহীহ মুসলিম ৮৩২-৮৩৩):}

\begin{quote}
উকবাহ ইবনে আমির (রাঃ) থেকে: নবীজি (সা.) \textbf{তিন সময়ে নামাজ পড়তে নিষেধ করেছেন} (নফল) — \\
১. \textbf{সূর্যোদয়ের সময়} — যতক্ষণ না সূর্য পূর্ণ উদিত হয় (একটি বর্শার উচ্চতা)। \\
২. \textbf{সূর্য ঠিক মাথার উপরে (দুপুরে)} — যতক্ষণ না ঢলে যায়। \\
৩. \textbf{সূর্যাস্তের সময়} — যতক্ষণ না সূর্য সম্পূর্ণ অস্ত যায়।
\end{quote}

\begin{quote}
\textbf{নোট:} এই নিষেধ \textbf{নফল নামাজের} জন্য — \textbf{ফরজ ওয়াক্ত-ভিত্তিক} নামাজ (যেমন যোহর ঢলে পড়লে, ফজর-যোহর আগের দিন মিস হলে) এ নিষেধে পড়ে না — সময় হলে ফরজ পড়া আবশ্যক। \textbf{কাযা নামাজ} নিষিদ্ধ সময়ে পড়া যায় কি না — মাযহাবভেদে (হানাফি: নিষিদ্ধ; জমহুর: মাকরূহ, তবে জরুরি হলে)।
\end{quote}

\textbf{অন্য নিষিদ্ধ সময়:}
\begin{itemize}
\item সূর্যাস্তের \textbf{লাল আভা} (শাফাক) — মাগরিব-আসরের সীমা; অস্ত-পরবর্তী লাল রঙে নফল নয়।
\item \textbf{জুমার খুতবার সময়} — খুতবা চলাকালে নফল নামাজ মাকরূহ (বুখারি ৯৩৪-এর আলোকে)।
\end{itemize}
\medskip\hrule\medskip

\subsection*{৪. সুতরা (সামনে রাখা) — আদব ও হুকুম}

\textbf{সহীহ হাদিস (সহীহ বুখারি ৪৯৬):}

\begin{quote}
নবীজি (সা.) বলেছেন — \textbf{\lat{“}সুতরা সামনে রাখা বান্দার উপর আবশ্যক — সুতরা সামনে না থাকলে, (কেউ পাশ দিয়ে গেলে) কুকুর কিংবা গাধা পাশ দিয়ে গেলেও (যেন সেটা থামে, কারণ সুতরা তার সামনে দিয়ে যাওয়া আটকায়)।\lat{”}} (আবু সাঈদ খুদরি — নবীজি (সা.) সুতরা ছাড়া নামাজ অপছন্দ করতেন)
\end{quote}

\textbf{সুতরার নিয়ম:}
\begin{itemize}
\item \textbf{মসজিদে:} স্তম্ভ, জায়নামাজের প্রান্ত, মুসহাফ (দাঁড়ানো), ব্যাগ — যে কোনো বস্তু। (কাতারে নামাজীরা পরস্পর সুতরা — ইমামের সামনে সুতরা হলে মুক্তাদি সুতরা নয় — জমহুর মত)।
\item \textbf{মাঠ-খোলা জায়গায়:} লাঠি, পাথর, কাপড়-রেখা।
\item \textbf{সুতরার দূরত্ব:} নামাজীর সামনে ১ কনুই থেকে ৩ হাত (প্রায় ০.৫-১.৫ মিটার)।
\item \textbf{সুতরার মাঝে কেউ গেলে:} সুতরার পেছন দিয়ে চলা জায়েয; সুতরার ভেতরে (নামাজী ও সুতরার মাঝে) গেলে নামাজ ভাঙে না, তবে আদব-লঙ্ঘন।
\end{itemize}
\medskip\hrule\medskip

\subsection*{৫. নামাজে অতিরিক্ত কিছু (বিতর্কিত)}

\begin{itemize}
\item \textbf{কুরআন দেখে (মুসহাফ ধরে) নামাজ:} জায়েয (সব মাযহাবের মতভেদে — শাফেয়ি-হাম্বলি জায়েয; হানাফি মাকরূহ; মালেকি জায়েয) — তবে \textbf{কুরআন পড়া স্মৃতি থেকে} হলে উত্তম। (বিতর্কিত — সাধারণ মুসল্লি মুসহাফ ধরে নফল-কুরআন-নামাজ করতে পারেন; ফরজ-এ স্মৃতি-পঠন উত্তম)
\item \textbf{মোবাইলে কুরআন দেখে নামাজ:} আধুনিক ফতোয়ায় অনুমোদন আছে (নির্বাচিত), তবে ফরজ-এ কুরআন-মুখস্থ পড়াই পূর্ণতা; মোবাইল-পড়া \textbf{নফল-তিলাওয়াত-নামাজে} বেশি ব্যবহৃত। — \textbf{সতর্কতা:} মোবাইল ফোন নামাজে বেজে ওঠা ও হাতে ধরা খুশু নষ্ট করে — প্রয়োজন না থাকলে পরিহার।
\end{itemize}
\medskip\hrule\medskip

\subsection*{৬. প্রশ্নোত্তর}

\begin{table}[h]\centering\small
\begin{tabular}{ll}
প্রশ্ন & উত্তর \\
\hline
ভুলে কথা বললে নামাজ? & জমহুর: ভাঙে (সাহু সিজদা নয়); মালেকি-হাম্বলি: সাহু সিজদায় ঠিক হতে পারে \\
খাবার খেলে? & ভাঙে (সব মাযহাব) \\
সূর্যোদয়-মধ্যাহ্ন-সূর্যাস্তে নফল? & নিষেধ — ফরজ পড়া জায়েয \\
মোবাইল বাজলে? & নামাজ ভাঙে না, মাকরূহ; শান্ত-রেখে নামাজ \\
সুতরা বাধ্যতামূলক? & সুন্নত/অত্যন্ত গুরুত্বপূর্ণ (বুখারি ৪৯৬-এর আলোকে) \\
\hline
\end{tabular}
\end{table}

\medskip\hrule\medskip

\subsection*{৭. রেফারেন্স}

\begin{itemize}
\item সহীহ বুখারি — ৩৭৪, ৪৯৬, ৭৮০, ৯৩৪, ১২১৭
\item সহীহ মুসলিম — ৫৩৭, ৫৬০, ৮৩২-৮৩৩
\item সুনানে আবু দাউদ — ৫৬৬-৫৬৯ (সুতরা ও নামাজ)
\item আল-ফিকহুল ইসলামি ওয়া আদিল্লাতুহু (আল-যুহাইলি) — ভঙ্গকারী ও নিষিদ্ধ সময়ের মাযহাব-সারণি
\end{itemize}
\begin{quote}
\textbf{দ্রষ্টব্য:} ভুলে ফরজ-ওয়াজিব ছুটে গেলে, মাসবুকের জটিল সংশোধন ও নামাজের ভুলের ধাপে ধাপে প্রতিকার — পার্ট ৭খ-এ।
\end{quote}

\end{document}
